\begin{tabular}{llll}
\toprule
variable & tipo & descripción & uso en el análisis \\
\midrule
mag & Numérica continua & Magnitud reportada del evento sísmico en el catálogo USGS. & Variable respuesta; se analiza solo dentro de la muestra truncada con mag >= 4. \\
depth & Numérica continua & Profundidad estimada del evento, usualmente en kilómetros según USGS. & Predictor geofísico interpretable; permite revisar si la magnitud reportada cambia con la profundidad del evento. \\
log\_depth\_plus\_1 & Numérica transformada & log(depth + 1), calculado como log1p(depth.clip(lower=0)). & Transformación usada porque depth es asimétrica y tiene valores extremos; estabiliza la relación lineal. \\
latitude & Numérica continua & Latitud del epicentro del evento. & Control espacial; captura diferencias promedio asociadas a la ubicación norte-sur. \\
longitude & Numérica continua & Longitud del epicentro del evento. & Control espacial; junto con latitude resume patrones geográficos del catálogo. \\
nst & Numérica discreta & Número de estaciones utilizadas para localizar o caracterizar el evento. & Variable técnica de cobertura: más estaciones pueden asociarse con eventos mejor registrados o de mayor alcance. \\
gap & Numérica continua & Brecha azimutal de estaciones alrededor del epicentro; valores mayores indican peor cobertura angular. & Variable técnica de geometría del registro; ayuda a controlar diferencias de calidad/cobertura. \\
dmin & Numérica continua & Distancia mínima entre el evento y una estación reportante. & Variable técnica de proximidad del registro; puede asociarse con precisión y disponibilidad de información. \\
rms & Numérica continua & Error cuadrático medio del ajuste del evento reportado por el catálogo. & Variable técnica de ajuste; se incluye como característica del registro, no como causa física directa. \\
month & Numérica discreta & Mes UTC derivado desde time. & Control temporal descriptivo para revisar estacionalidad o cambios de composición del catálogo. \\
hour & Numérica discreta & Hora UTC derivada desde time. & Control temporal descriptivo; no es predictor físico central. \\
magType & Categórica & Tipo o escala de magnitud reportada, por ejemplo mb, mww u otras codificaciones. & Permite comparar diferencias promedio por método de medición/codificación. \\
magType\_grouped & Categórica agrupada & Versión agrupada de magType que conserva niveles frecuentes y reúne niveles escasos como Other. & Reduce dummies con pocas observaciones y mejora estabilidad del modelo lineal. \\
net & Categórica & Red u organismo que reporta el evento. & Controla diferencias promedio asociadas a la institución o red de reporte. \\
net\_grouped & Categórica agrupada & Versión agrupada de net que conserva redes frecuentes y reúne niveles escasos como Other. & Evita categorías con muy poca información en los modelos con indicadores. \\
region\_from\_place & Categórica derivada & Región aproximada extraída desde el texto place. & Resumen geográfico interpretable para EDA; no reemplaza latitude-longitude. \\
region\_grouped & Categórica agrupada & Versión agrupada de region\_from\_place. & Permite describir regiones frecuentes sin saturar tablas o gráficos con niveles raros. \\
place & Texto & Descripción textual del lugar del evento. & Contexto y construcción de region\_from\_place. \\
time & Fecha-hora & Fecha y hora UTC del evento. & Filtro del año 2024 y creación de variables temporales. \\
id & Identificador & Identificador único del evento en el catálogo USGS. & Control de duplicados. \\
\bottomrule
\end{tabular}
